% =====================================================================
% Resumen (ES) y abstract (EN). Van antes del índice, sin numerar.
% Las cifras citadas aquí proceden del Study \studyid y deben coincidir
% con las de los capítulos 5 a 7.
% =====================================================================
\chapter*{Resumen}
\addcontentsline{toc}{chapter}{Resumen}

Este trabajo evalúa si un sistema de cinco agentes especializados, combinados por un meta-agente que
aprende sus pesos, es capaz de aprender una ordenación transversal de acciones que conserve
capacidad predictiva fuera de muestra. La pregunta se plantea deliberadamente en términos de
predicción y no de rentabilidad retrospectiva: el objetivo no es encontrar la cartera más rentable
de la historia conocida, sino medir si existe señal y cuánta de esa señal sobrevive a la
implementación.

El universo es la composición histórica del S\&P 500 sobre un panel \textit{point-in-time}
construido a partir de fechas reales de publicación, de modo que ninguna decisión utiliza
información no disponible en su fecha. Sobre un catálogo cerrado de variables se ejecuta un único
\textit{Model Study} que optimiza de forma secuencial por Rank-IC, dejando 2025--2026 como era
reservada que no interviene en ninguna decisión. La cartera y los ocho perfiles de inversión se
construyen después de congelar el ganador, y nunca influyen en su selección.

El sistema alcanza un Rank-IC medio de 0,1004 con IC-IR 0,744 y 71,79\,\% de cohortes positivas en
la ventana de selección, frente a 0,0130 del mejor \textit{baseline} determinista. La ponderación
aprendida por el meta-agente supera en un 52\,\% a la ponderación uniforme. Siete de ocho contrastes
de robustez respaldan la existencia de señal --- permutación con $p=0{,}0001$, placebos de etiqueta,
\textit{bootstrap} por bloques, exclusión de eras, estabilidad entre semillas, carteras aleatorias
de riesgo emparejado y neutralización de estilo ---, mientras que el Deflated Sharpe de 0,930 no
alcanza su umbral y se reporta como tal. En la era reservada la cartera obtiene un exceso geométrico
del 14,21\,\% con IR 0,959, si bien su Rank-IC cerrado todavía carece de potencia estadística.

El resultado central para la práctica es el coeficiente de transferencia, 0,247: la cartera
materializa aproximadamente una cuarta parte de la señal disponible. El cuello de botella no es el
modelo sino la construcción de cartera, una conclusión que el barrido diagnóstico de veintitrés
variantes confirma de forma controlada.

\vspace{1em}
\noindent\textbf{Palabras clave:} aprendizaje automático, selección de acciones, \textit{point-in-time},
Rank-IC, sistemas multiagente, sobreajuste de \textit{backtest}.

\chapter*{Abstract}
\addcontentsline{toc}{chapter}{Abstract}

This work evaluates whether a system of five specialised agents, combined by a meta-agent that
learns their weights, can learn a cross-sectional ranking of stocks that retains out-of-sample
predictive power. The question is deliberately framed in terms of prediction rather than
retrospective profitability: the goal is not to find the most profitable portfolio in known history,
but to measure whether signal exists and how much of it survives implementation.

The universe is the historical composition of the S\&P 500 over a point-in-time panel built from
actual publication dates, so that no decision uses information unavailable at its date. A single
Model Study is run over a closed catalogue of variables, optimising sequentially by Rank-IC and
holding out 2025--2026 as a reserved era that takes part in no decision. The portfolio and the eight
investor profiles are constructed only after the winner is frozen, and never influence its
selection.

The system attains a mean Rank-IC of 0.1004 with an IC-IR of 0.744 and 71.79\,\% positive cohorts in
the selection window, against 0.0130 for the best deterministic baseline. The meta-agent's learned
weighting outperforms uniform weighting by 52\,\%. Seven of eight robustness tests support the
existence of signal --- permutation with $p=0.0001$, label placebos, block bootstrap, era exclusion,
seed stability, risk-matched random portfolios and style neutralisation --- while the Deflated
Sharpe of 0.930 falls short of its threshold and is reported as such. In the reserved era the
portfolio achieves a geometric excess return of 14.21\,\% with an IR of 0.959, although its closed
Rank-IC still lacks statistical power.

The central practical finding is the transfer coefficient of 0.247: the portfolio materialises
roughly one quarter of the available signal. The bottleneck is not the model but portfolio
construction, a conclusion confirmed under controlled conditions by a diagnostic sweep of
twenty-three variants.

\vspace{1em}
\noindent\textbf{Keywords:} machine learning, stock selection, point-in-time data, Rank-IC,
multi-agent systems, backtest overfitting.
